```latex
\documentclass[11pt]{article}
\usepackage[margin=1in]{geometry}
\usepackage{graphicx}
\usepackage{amsmath,amssymb}
\usepackage{booktabs}
\usepackage{hyperref}
\usepackage{authblk}
\usepackage{caption}
\usepackage{xcolor}

\title{Autoscientist: A Multi-Agent Framework for Data Efficiency in Medical Imaging}
\author{Autoscientist Collaboration}
\date{}

\begin{document}
\maketitle

\begin{abstract}
Acquiring large labeled medical imaging datasets is expensive, motivating principled study of whether \emph{matching} training data to a target domain can improve sample efficiency, particularly at small training sizes. We present \textsc{Autoscientist}, a multi-agent framework that autonomously designs, executes, and interprets controlled deep learning experiments in medical imaging. Using \textsc{Autoscientist}, we test the hypothesis that domain-matched training data yields disproportionately larger gains over unmatched data at smaller sample sizes. We train chest radiograph classifiers on NIH ChestX-ray14 subsets of size $N \in \{1\text{k}, 5\text{k}, 25\text{k}, 100\text{k}\}$ under matched and unmatched conditions and evaluate transfer to the PAD-CHEST test set. A permutation-based interaction test fails to reject the null ($F=0.56$, $p=0.217$). Only $N=5\text{k}$ shows a significant matched-vs-unmatched difference after correction (adjusted $p=0.021$), while at $N=25\text{k}$ the direction unexpectedly reverses (matched 0.446 vs.\ unmatched 0.477; not significant). All AUROC values cluster near chance (0.43--0.49), indicating substantial cross-domain generalization failure that confounds the sample-efficiency signal. We discuss what these negative results imply for designing data-efficient medical imaging pipelines and how an autonomous agentic workflow surfaces such confounds early.
\end{abstract}

\section{Introduction}

Medical imaging models are notoriously data-hungry, and label acquisition is constrained by clinician time, privacy regulation, and institutional heterogeneity \cite{rajpurkar2017chexnet,cohen2020limits}. A common assumption in the data-efficiency literature is that distributional \emph{matching} between training and target domains is most beneficial when training data is scarce: at small $N$, every sample matters, and reducing covariate shift should yield outsized gains. Whether this assumption holds empirically for modern chest radiograph classifiers is unclear, with prior cross-institutional studies reporting mixed and often distribution-shift-dominated results \cite{zech2018variable}.

Manually testing such hypotheses is laborious: it requires careful curation of training subsets, repeated training under matched random seeds, and statistically principled comparison across conditions. We present \textsc{Autoscientist}, a multi-agent framework in which specialized agents handle hypothesis formalization, data curation, training orchestration, evaluation, and interpretation. The framework's central design goal is to make medical imaging experiments \emph{reproducible by construction} and \emph{statistically honest by default}.

In this paper, we use \textsc{Autoscientist} to ask: \emph{Does matching training data to a target domain improve sample efficiency more at smaller training sizes?} Our contributions are: (1) a description of the multi-agent workflow; (2) a controlled $2 \times 4$ experiment (matching condition $\times$ training size) with permutation-based interaction testing; and (3) a transparent discussion of negative and near-null results, including a counterintuitive reversal at $N=25\text{k}$ and evidence that cross-domain transfer failure dominates the observed signal.

\section{Methods}

\subsection{Multi-Agent Framework}

\textsc{Autoscientist} comprises five cooperating agents: (i) a \emph{Hypothesis Agent} that formalizes the scientific question and pre-registers the analysis plan; (ii) a \emph{Data Agent} that constructs matched and unmatched training subsets and target test sets; (iii) a \emph{Training Agent} that orchestrates model training across conditions and random seeds; (iv) an \emph{Evaluation Agent} that computes discrimination (AUROC) and calibration (ECE, Brier) metrics; and (v) an \emph{Interpretation Agent} that performs statistical inference, multiple-testing correction, and figure generation. Each agent communicates via structured handoffs (JSON-typed artifacts), and all randomness, seeds, and dataset partitions are logged.

\subsection{Datasets and Matching}

Training data is drawn from the NIH ChestX-ray14 corpus; the held-out generalization test set is drawn from PAD-CHEST. In the \emph{matched} condition, training subsets are resampled to approximate the marginal distribution of age, sex, and view position observed in the PAD-CHEST test cohort. In the \emph{unmatched} condition, training subsets are sampled uniformly at random from NIH. Both conditions are evaluated at $N \in \{1{,}000;\ 5{,}000;\ 25{,}000;\ 100{,}000\}$.

\subsection{Models and Training}

For each $(N, \text{condition})$ cell we train a standard convolutional backbone (ImageNet-initialized) for binary disease classification with cross-entropy loss, fixed optimizer, and a held-out NIH validation split for early stopping. Multiple random seeds are used per cell. Class imbalance is preserved as in the source data, which at the smallest $N$ yields very few positive cases (e.g., 8 positives at $N=1{,}000$).

\subsection{Statistical Analysis}

The primary test is a permutation test for the interaction between matching condition and $\log N$ on test AUROC. Pairwise comparisons within each $N$ use two-sample $t$-tests with Benjamini--Hochberg correction across the four sizes. Calibration differences (ECE, Brier) are tested analogously. All tests, multiple-testing thresholds, and reporting precision were specified by the Hypothesis Agent prior to result inspection.

\section{Results}

\subsection{No Significant Interaction Between Matching and Sample Size}

The pre-registered permutation test for the matching $\times$ $\log N$ interaction was not significant ($F = 0.558$, interaction coefficient $= 0.0085$, permutation $p = 0.217$). The central hypothesis --- that matched data provides larger gains at smaller $N$ --- is not supported.

\subsection{AUROC by Condition}

Pairwise AUROC results are summarized in Table~\ref{tab:auroc}. The only comparison surviving multiple-testing correction is at $N=5{,}000$, where matched data outperforms unmatched (0.491 vs.\ 0.425; adjusted $p = 0.021$). At $N=1{,}000$ and $N=100{,}000$ matched data trends higher but does not reach significance. At $N=25{,}000$, the direction unexpectedly reverses: unmatched data yields a higher mean AUROC (0.477) than matched (0.446), though this difference is not significant (adjusted $p = 0.669$).

\begin{table}[h]
\centering
\caption{Pairwise AUROC results across training sizes. Adjusted $p$-values use Benjamini--Hochberg correction.}
\label{tab:auroc}
\begin{tabular}{rcccc}
\toprule
$N$ & Matched & Unmatched & $\Delta$ (M$-$U) & Adj. $p$ \\
\midrule
1{,}000   & 0.486 & 0.427 & $+0.059$ & 0.669 \\
5{,}000   & 0.491 & 0.425 & $+0.066$ & \textbf{0.021} \\
25{,}000  & 0.446 & 0.477 & $-0.031$ & 0.669 \\
100{,}000 & 0.465 & 0.432 & $+0.033$ & 0.669 \\
\bottomrule
\end{tabular}
\end{table}

Critically, \emph{all} AUROC means lie between 0.43 and 0.49, i.e., at or below chance, indicating substantial cross-domain generalization failure from NIH to PAD-CHEST. This near-random regime makes it difficult to attribute remaining variation to matching versus noise.

\subsection{Calibration}

Calibration results are heterogeneous and noisy (Table~\ref{tab:cal}). At $N=1{,}000$ and $N=25{,}000$, matched models are markedly \emph{worse} calibrated than unmatched (e.g., ECE 0.313 vs.\ 0.097 at $N=1\text{k}$). At $N=5{,}000$ and $N=100{,}000$, matched models are slightly better calibrated. No calibration difference survives multiple-testing correction (all adjusted $p = 1.0$).

\begin{table}[h]
\centering
\caption{Calibration metrics by condition (ECE = Expected Calibration Error).}
\label{tab:cal}
\begin{tabular}{rcccc}
\toprule
$N$ & Matched ECE & Unmatched ECE & Matched Brier & Unmatched Brier \\
\midrule
1{,}000   & 0.313 & 0.097 & 0.210 & 0.067 \\
5{,}000   & 0.088 & 0.118 & 0.060 & 0.073 \\
25{,}000  & 0.327 & 0.134 & 0.248 & 0.111 \\
100{,}000 & 0.048 & 0.067 & 0.048 & 0.063 \\
\bottomrule
\end{tabular}
\end{table}

\section{Discussion}

Our pre-registered analysis fails to support the hypothesis that matched training data improves sample efficiency more at smaller $N$. Three findings are worth highlighting.

\paragraph{The $N=5{,}000$ exception.} The only significant matched-over-unmatched advantage occurs at $N=5{,}000$, which also coincides with the best calibration for matched models. This may reflect a narrow regime where the model has enough data to exploit distributional alignment but not enough to overcome shift on its own. As an isolated result among four comparisons, however, this is at best suggestive.

\paragraph{The $N=25{,}000$ reversal.} The direction flip at $N=25{,}000$ is striking, although not statistically significant. Plausible explanations include (i) a transitional regime where the model shifts from data-limited to capacity-limited, and matching-induced distributional narrowing reduces effective diversity; (ii) artifacts of the matching procedure itself at medium scale; and (iii) seed-level stochasticity amplified by the near-chance AUROC landscape. Disentangling these would require denser sampling of $N$ and more seeds per cell --- both natural next steps for \textsc{Autoscientist} to schedule autonomously.

\paragraph{Cross-domain transfer dominates the signal.} All AUROCs cluster near 0.5, indicating that the NIH$\rightarrow$PAD-CHEST shift is the dominant effect, swamping any matching-driven sample-efficiency benefit. In this regime, conclusions about ``sample efficiency'' are themselves suspect: we are largely measuring generalization failure. Future iterations should either (i) test on an in-distribution NIH hold-out to isolate sample-efficiency effects, or (ii) use stronger matching variables (e.g., scanner, hospital) and domain-adaptive training objectives.

\paragraph{Implications for autonomous experimentation.} A central value of the \textsc{Autoscientist} framework is precisely its willingness to report negative, mixed, and confounded results. The Interpretation Agent flagged the near-chance AUROCs and the direction reversal as red